\chapter{Strong Divisibility Sequences}

\section{Lucas Sequences}

The French mathematician \'Edouard Lucas observed that the familiar Fibonacci sequence was just a single instance of a larger family of recursively defined sequences that all share similar properties. 

\begin{definition}
  Consider two integers $a$ and $b$. The \emph{Lucas sequence} $u_n$ at $a$ and $b$ is defined by
  \begin{equation*}
    u_0:= 0, \qquad u_1:=1,\qquad\text{and}\qquad u_{n+2}:= au_{n+1}- bu_n.
  \end{equation*}
  If the parameters $a$ and $b$ are not clear from the context, we will write $u(a,b)_n$ for $u_n$. When $a$ and $b$ are relatively prime, we say that the Lucas sequence $(u_n)$ is \emph{reduced}.
\end{definition}

\begin{example}
  Some special cases of Lucas sequences are famous:
  \begin{enumerate}
  \item When $a=1$ and $b=-1$, then the sequence $(u_n)$ is the Fibonacci sequence
    \begin{equation*}
      0,1,1,2,3,5,8,13,21,34,55,89,\ldots
    \end{equation*}
    When we look at the Fibonacci sequence modulo a prime $p$, we see clear patterns emerging:
    \begin{equation*}
      \newcolumntype{C}{>{\raggedleft\arraybackslash$}m{.8em}<{$}}
      \renewcommand{\arraystretch}{1.25}
      \begin{tabular}{l|CCCCCCCCCCCCCCCC}
        $n$             & 0 & 1 & 2 & 3 & 4 & 5 & 6 & 7 & 8 & 9 & 10 & 11 & 12 & 13 & 14 & 15 \\ \midrule
        $F_n\pmod{2}$ & 0 & 1 & 1 & 0 & 1 & 1 & 0 & 1 & 1 & 0 & 1 & 1 & 0 & 1 & 1 & 0 \\
        $F_n\pmod{3}$ & 0 & 1 & 1 & -1 & 0 & -1 & -1 & 1 & 0 & 1 & 1 & -1 & 0 & -1 & -1 & 1 \\
        $F_n\pmod{5}$ & 0 & 1 & 1 & 2 & -2 & 0 & -2 & -2 & 1 & -1 & 0 & -1 & -1 & -2 & 2 & 0 \\
        $F_n\pmod{7}$ & 0 & 1 & 1 & 2 & 3 & -2 & 1 & -1 & 0 & -1 & -1 & -2 & -3 & 2 & -1 & 1
      \end{tabular}
    \end{equation*}
  \item When $a=2$ and $b=-1$, then the sequence $(u_n)$ is the \emph{Pell sequence}
    \begin{equation*}
      0,1,2,5,12,29,70,169,408,985,\ldots
    \end{equation*}
    When we look at the Pell sequence modulo a prime $p$, we again see clear patterns emerging:
    \begin{equation*}
      \newcolumntype{C}{>{\raggedleft\arraybackslash$}m{.8em}<{$}}
      \renewcommand{\arraystretch}{1.25}
      \begin{tabular}{l|CCCCCCCCCCCCCCCC}
        $n$             & 0 & 1 & 2 & 3 & 4 & 5 & 6 & 7 & 8 & 9 & 10 & 11 & 12 & 13 & 14 & 15 \\ \midrule
        $P_n\pmod{2}$ & 0 & 1 & 0 & 1 & 0 & 1 & 0 & 1 & 0 & 1 & 0 & 1 & 0 & 1 & 0 & 1 \\
        $P_n\pmod{3}$ & 0 & 1 & -1 & -1 & 0 & -1 & 1 & 1 & 0 & 1 & -1 & -1 & 0 & -1 & 1 & 1 \\
        $P_n\pmod{5}$ & 0 & 1 & 2 & 0 & 2 & -1 & 0 & -1 & -2 & 0 & -2 & 1 & 0 & 1 & 2 & 0 \\
        $P_n\pmod{7}$ & 0 & 1 & 2 & -2 & -2 & 1 & 0 & 1 & 2 & -2 & -2 & 1 & 0 & 1 & 2 & -2
      \end{tabular}
    \end{equation*}
  \item In \cref{exercise:mersenne-lucas-sequence} we saw that the sequence $(M_n)$ of Mersenne numbers satisfies the recurrence
    \begin{equation*}
      M_{n+2}=3M_{n+1}-2M_n.
    \end{equation*}
    Thus, the Mersenne numbers also form an example of a Lucas sequence. When we look at the Mersenne numbers modulo a prime $p$, we obtain the following:
    \begin{equation*}
      \newcolumntype{C}{>{\raggedleft\arraybackslash$}m{.8em}<{$}}
      \renewcommand{\arraystretch}{1.25}
      \begin{tabular}{l|CCCCCCCCCCCCCCCC}
        $n$             & 0 & 1 & 2 & 3 & 4 & 5 & 6 & 7 & 8 & 9 & 10 & 11 & 12 & 13 & 14 & 15 \\ \midrule
        $M_n\pmod{2}$ & 0 & 1 & 1 & 1 & 1 & 1 & 1 & 1 & 1 & 1 & 1 & 1 & 1 & 1 & 1 & 1 \\
        $M_n\pmod{3}$ & 0 & 1 & 0 & 1 & 0 & 1 & 0 & 1 & 0 & 1 & 0 & 1 & 0 & 1 & 0 & 1 \\
        $M_n\pmod{5}$ & 0 & 1 & -2 & 2 & 0 & 1 & -2 & 2 & 0 & 1 & -2 & 2 & 0 & 1 & -2 & 2 \\
        $M_n\pmod{7}$ & 0 & 1 & 3 & 0 & 1 & 3 & 0 & 1 & 3 & 0 & 1 & 3 & 0 & 1 & 3 & 0
      \end{tabular}
    \end{equation*}
  \item More generally, for any two distinct integers $a$ and $b$, the sequence defined by
    \begin{equation*}
      u_n := \frac{a^n-b^b}{a-b}
    \end{equation*}
    is a Lucas sequence satisfying
    \begin{equation*}
      u_0=0,\qquad u_1=1,\qquad\text{and}\qquad u_{n+2}=(a+b)u_{n+1}+abu_n. 
    \end{equation*}
  \end{enumerate}
\end{example}

\begin{theorem}[Binet's Formula]\label{theorem:binets-formula}
  Consider two $a$ and $b$ such that $a^2-4b\neq 0$. Suppose $\alpha$ and $\beta$ are roots of the polynomial
  \begin{equation*}
    x^2-ax+b=0.
  \end{equation*}
  Then we have $\alpha\neq \beta$, and the Lucas sequence $(u_n)$ at $a$ and $b$ satisfies
  \begin{equation*}
    u_n=\frac{\alpha^n-\beta^n}{\alpha-\beta}.
  \end{equation*}
\end{theorem}

\begin{proof}
  Given that $\alpha$ and $\beta$ are roots of the polynomial $x^2-ax+b$, we have $\alpha^2=a\alpha-b$ and $\beta^2=a\beta-b$. This implies the identities
  \begin{equation*}
    \alpha^{n+2}=a\alpha^{n+1}-b\alpha^n\qquad\text{and}\qquad \beta^{n+2}=a\beta^{n+1}-b\beta^n.
  \end{equation*}
  We will now prove Binet's formula by induction on $n$. There are two base cases, which are both immediate, since $\alpha^0-\beta^0=0$ and $\alpha^1-\beta^1=\alpha-\beta$. For the inductive step, we have
  \begin{equation*}
    \frac{\alpha^{n+2}-\beta^{n+2}}{\alpha-\beta}=\frac{a(\alpha^{n+1}-\beta^{n+1})-b(\alpha^n-\beta^n)}{\alpha-\beta}=au_{n+1}-bu_n=u_{n+2}.\qedhere
  \end{equation*}
\end{proof}

Binet's formula can be used to show that if $m\mid n$, then $u_m\mid u_n$. Sequences with this property are called \emph{divisibility sequences}. We will additionally assume that divisibility sequences $(a_n)$ satisfy $a_0=0$ and $a_1=0$. Such divisibility sequences are sometimes called \emph{normalized}.

\begin{definition}
  An integer sequence $(a_n)$ is said to be a \emph{divisibility sequence} if $a_0=0$, $a_1=1$, and it preserves the divisibility relation, meaning that it satisfies the condition
  \begin{equation*}
    m\mid n \Rightarrow a_m\mid a_n
  \end{equation*}
  for all $m$ and $n$.
\end{definition}

\begin{corollary}
  Consider two integers $a$ and $b$ such that $a^2-4b\neq 0$. Then the Lucas sequence $(u_n)$ at $a$ and $b$ is a divisibility sequence.
\end{corollary}

\begin{proof}
  Suppose that $n=km$. Then we have
  \begin{equation*}
    u_n=\frac{\alpha^n-\beta^n}{\alpha-\beta}=\frac{(\alpha^m-\beta^m)\sum_{i=0}^k \alpha^i\beta^{k-i}}{\alpha-\beta}=u_m\sum_{i=0}^k \alpha^i\beta^{k-i}
  \end{equation*}
  by the formula for the difference of $m$th powers. This shows that if $m\mid n$, then $u_m\mid u_n$.
\end{proof}

\begin{definition}
  A divisibility sequence is said to be a \emph{strong divisibility sequence} if $f$ preserves greatest common divisors, meaning that it satisfies the equation
  \begin{equation*}
    f(\gcd(m,n))=\gcd(f(m),f(n))
  \end{equation*}
  for all $m$ and $n$.
\end{definition}

\section{The Rank of Apparition}

\begin{definition}
  Consider an integer sequence $(a_n)$ and a positive integer $m$. The \emph{rank of apparition} $\rho_a(m)$ of $m$ is the least positive integer $n$ such that $m\mid a_n$.
\end{definition}

\begin{theorem}[Rank of Apparition Theorem]\label{theorem:galois-connections-divisibility}
  Suppose that $f:\mathbb{N}\to\mathbb{N}$ is a divisibility sequence. Then the following are equivalent:
  \begin{enumerate}
  \item There is a function $g:\mathbb{N}\to\mathbb{N}$ satisfying the logical equivalence
  \begin{equation*}
    g(m)\mid n \Leftrightarrow m\mid f(n).
  \end{equation*}
  \item The divisibility sequence $f$ is a strong divisibility sequence.
  \end{enumerate}
  If the first condition holds, then the function $g$ is equal to the rank of apparition function; that is, $g(m)=\rho_f(m)$ for every natural number $m$.
\end{theorem}

\begin{proof}
  Suppose first that there exists a function $g$ satisfying the logical equivalence
  \begin{equation*}
    g(m)\mid n \Leftrightarrow m\mid f(n).
  \end{equation*}
  Then it follows that
  \begin{align*}
    d\mid f(\gcd(m,n)) & \Leftrightarrow g(d)\mid \gcd(m,n) \\
                       & \Leftrightarrow g(d)\mid m \land g(d)\mid n \\
                       & \Leftrightarrow d\mid f(m) \land d\mid f(n) \\
                       & \Leftrightarrow d\mid \gcd(f(m),f(n)).
  \end{align*}
  Since $f(\gcd(m,n))$ and $\gcd(f(m),f(n))$ have the exact same set of divisors, they must be equal.
  
  For the converse, suppose that $f(0)=0$ and that $f$ preserves greatest common divisors. More generally, we will have for any finite list of natural numbers
  \begin{equation*}
    f(\gcd_i(n_i))=\gcd_i(f(n_i)).
  \end{equation*}
  Here, the expression $gcd_i(n_i)$ means that the greatest common divisor is taken over all the numbers $n_i$ in the list. Now consider for any natural number $m$ the ideal
  \begin{equation*}
    I_m := \left\{ k_1n_1+\cdots+ k_rn_r \mathbin{\big|} m\mid f(n_i) \text{ for all } i \right\}
  \end{equation*}
  and define the function $g$ by letting $g(m)$ be the unique natural number such that $I_m=(g(m))$.

  To see that $g$ is a divisibility sequence, suppose that $m\mid m'$. Then any linear combination of natural numbers $n$ such that $m'\mid f(n)$ is also a linear combination of natural numbers $n$ such that $m\mid f(n)$. In particular $g(m')$ is in the ideal $(g(m))$, which means that $g(m)\mid g(m')$.

  We claim that the function $g$ satisfies the logical equivalence
  \begin{equation*}
    g(m)\mid n\Leftrightarrow m\mid f(n).
  \end{equation*}
  If $m\mid f(n)$, then it follows that $n\in I_m$ so that $g(m)\mid n$. This proves the reverse direction.

  For the forward direction, suppose that $g(m)\mid n$. Then it also follows that $n\in I_m$, so that we can write
  \begin{equation*}
    n=\sum_i k_in_i
  \end{equation*}
  for some finite list of natural numbers $n_i$ satisfying $m\mid f(n_i)$. By the assumption that $m\mid f(n_i)$, it follows that $m\mid\gcd_i(f(n_i))$. Furthermore, since $\gcd_i(n_i)\mid n$ it follows that
  \begin{equation*}
    f(\gcd\nolimits_i(n_i))\mid f(n).
  \end{equation*}
  Combined with the fact that $m\mid\gcd_i(f(n_i))$ and the assumption that $f$ preserves greatest common divisors, we conclude that
  \begin{equation*}
    m\mid f(n).\qedhere
  \end{equation*}
\end{proof}

By the \crefnamed{theorem:order-theorem}{Order Theorem}, it follows that the function $g$ defined by $g(m):=\ord_m(a)$ satisfies the logical equivalence in the first condition of the \crefnamed{theorem:rank-of-apparition-theorem}{Rank of Apparition Theorem}. Thus we obtain the following corollary:

\begin{corollary}
  The function $n\mapsto a^n-1$ is a strong divisibility sequence. Thus, for any integer $a>1$ and any two natural numbers $m$ and $n$ we have
  \begin{equation*}
    \gcd(a^m-1,a^n-1)=a^{\gcd(m,n)}-1.
  \end{equation*}
\end{corollary}

\begin{proof}
  Note that the function $k\mapsto a^k-1$ is a divisibility sequence, since for $qm=n$ we have the identity
  \begin{equation*}
    a^n-1=(a^m)^q-1^q=(a^m-1)\sum_{k=0}^{q-1}a^{km}
  \end{equation*}
  by the formula for the difference of $q$th powers (\cref{exercise:difference-of-powers}). The fact that it is a strong divisibility sequence now follows directly from \cref{theorem:galois-connections-divisibility} and the \crefnamed{theorem:order-theorem}{Order Theorem}.
\end{proof}

\begin{remark}
  More generally, if $(P,\leq_P)$ and $(Q,\leq_Q)$ are two posets, a \emph{Galois connection} between $P$ and $Q$ consists of order preserving maps $f:P\to Q$ and $g:Q\to P$ satisfying the logical equivalence
  \begin{equation*}
    g(x)\leq_P y \Leftrightarrow x\leq_Q f(y)
  \end{equation*}
  for any $x\in Q$ and $y\in P$. In this case, the map $f$ is said to be the \emph{upper adjoint}, and the map $g$ is said to be the \emph{lower adjoint} of the Galois connection. Thus we have proven in this section that the function $n\mapsto a^n-1$ is the upper adjoint of a Galois connection, and we have used this fact to show that $f$ preserves greatest common divisors. \cref{theorem:galois-connections-divisibility} generalizes to the setting of arbitrary posets, and indeed, some readers may recognize it as a special case of the \emph{Adjoint Functor Theorem}.

  Another famous Galois connection on the natural numbers ordered by divisibility is the Fibonacci sequence. We will give the full statement of this fact in the exercises.
\end{remark}

\section{Primitive Sequences}

\begin{definition}
  Consider a divisibility sequence $n\mapsto a_n$ indexed by the natural numbers, with $a_1=1$. The \emph{primitive sequence} of $a_n$ is the sequence $(u_n)$ of rational numbers defined by
  \begin{equation*}
    u_n:=\prod_{d\mid n}a_d^{\mu(n/d)}.
  \end{equation*}
\end{definition}

\begin{theorem}
  The primitive sequence $(u_n)$ of a divisibility sequence $(a_n)$ satisfies the identity
  \begin{equation*}
    a_n=\prod_{d\mid n} u_d.
  \end{equation*}
\end{theorem}

\begin{proof}
  By interchanging the product, we find that
  \begin{equation*}
    \prod_{e\mid n}\prod_{d\mid e} a_{d}^{\mu(e/d)}=\prod_{d\mid n}\prod_{c\mid\frac{n}{d}} a_{d}^{\mu(c)}=\prod_{d\mid n}a_d^{\sum_{c\mid\frac{n}{d}}\mu(c)}=\prod_{d\mid n}a_d^{\varepsilon(n/d)}=a_n. \qedhere
  \end{equation*}
\end{proof}

\section{LCM-sequences}

\begin{definition}
  The \emph{lcm-sequence} $(m_n)$ of an integer sequence $(a_n)$ is given by $m_1:=1$ and
  \begin{equation*}
    m_n := \frac{\lcm(a_1,\ldots,a_n)}{\lcm(a_1,\ldots,a_{n-1})}\qquad\text{for $n\geq 2$}.
  \end{equation*}
\end{definition}

\begin{theorem}
  Consider a divisibility sequence $(a_n)$ indexed by the natural numbers, with $a_1=1$. Then the following are equivalent:
  \begin{enumerate}
  \item The sequence $(a_n)$ is a strong divisibility sequence; that is,
    \begin{equation*}
      a_{\gcd(m,n)}=\gcd(a_m,a_n)
    \end{equation*}
    for every $m$ and $n$.
  \item \textup{\cite{ward1936}} For each prime $p$, we have $p^k\mid a_n$ if and only if $\rho_k(p)\mid n$, where $\rho_k(p)$ is the rank of apparition of $p^k$.
  \item Each $u_n$ in primitive sequence $(u_n)$ of $(a_n)$ is an integer, and furthermore $u_m$ and $u_n$ are relatively prime for distinct $m$ and $n$.
  \item \textup{\cite{nowicki2015}} The sequence $(a_n)$ satisfies the following divisor-product formula, where $(m_n)$ is the lcm-sequence of $(a_n)$:
    \begin{equation*}
      a_n = \prod_{d\mid n}m_d.
    \end{equation*}
  \end{enumerate}
\end{theorem}

The following corollary, due to Morgan Ward \cite{ward1936}, is a generalization of a theorem of \'Edouard Lucas \cite{lucas1878}.

\begin{corollary}
  If $(a_n)$ is a strong divisibility sequence, then
  \begin{equation*}
    \prod_{i=1}^ka_i \Bigm| \prod_{i={n+1}}^{n+k}a_i.
  \end{equation*}
\end{corollary}

\section{The Lifting the Exponent Lemma}

For the following lemma, recall that $\nu_p(a-b)=m$ if and only if $a\equiv b\pmod{p^m}$ but $a\nequiv b\pmod{p^{m+1}}$. We also recall that for any prime $p$ and any $0<i<p$, the binomial coefficient $\binom{p}{i}$ is divisible by $p$. 

\begin{lemma}\label{lemma:lte-step}
  Consider an odd prime $p$, and two integers $a\equiv b\pmod{p}$ such that neither $a$ nor $b$ is divisible by $p$. Then
  \begin{equation*}
    \nu_p(a^p-b^p)=\nu_p(a-b)+1.
  \end{equation*}
\end{lemma}

\begin{proof}
  Write $s:=\nu_p(a-b)$, so that $s\geq 1$ by the assumption that $a\equiv b\pmod{p}$. Then we have $a\equiv b\pmod{p^s}$, so we can write $b=a+tp^s$, where $p\nmid t$. The value of $b^p$ can therefore be computed using the \crefnamed{theorem:binomial-theorem}{Binomial Theorem} by
  \begin{equation*}
    b^p=(a+tp^s)^p=\sum_{i=0}^p\binom{p}{i}a^i(tp^s)^{p-i} \equiv a^p+a^{p-1}tp^{s+1}\pmod{p^{s+2}}.
  \end{equation*}
  Since neither $a$ nor $t$ is divisible by $p$, it follows that $a^p\nequiv b^p\pmod{p^{s+2}}$. On the other hand, we see from the binomial expansion of $b^p$ that $a^p\equiv b^p\pmod{p^{s+1}}$. Therefore we conclude that
  \begin{equation*}
    \nu_p(a^p-b^p)=s+1.\qedhere
  \end{equation*}
\end{proof}

The previous lemma can be iterated to obtain a general lemma called the \emph{Lifting the Exponent Lemma}. This lemma was popularized by the Romanian mathematician Mihai Manea \cite{manea2006}.

\begin{lemma}[Lifting the Exponent]\label{lemma:lte}
  Consider an odd prime $p$, and let $x\equiv y\pmod{p}$, where neither $x$ nor $y$ is divisible by $p$. Then for any natural number $n\geq 1$, we have
  \begin{equation*}
    \nu_p(x^n-y^n)=\nu_p(x-y)+\nu_p(n).
  \end{equation*}
\end{lemma}

\begin{proof}
  Write $n=p^km$ with $p\nmid m$, so that $\nu_p(n)=k$, and write $s=\nu_p(x-y)$. Since we have assumed that $x\equiv y\pmod{p}$, we have $s\geq 1$. We will prove the claim by induction on $k$.

  In the base case, where $k=0$, the claim is simply that
  \begin{equation*}
    \nu_p(x^m-y^m)=s.
  \end{equation*}
  To see this, write $y=x+tp^s$, where $p\nmid t$. Then we have
  \begin{equation*}
    y^m=(x+tp^s)^m\equiv x^m+mx^{m-1}tp^s \pmod{p^{s+1}}.
  \end{equation*}
  Now observe that $m$, $x$, and $t$ are not divisible by $p$ by assumption, so it follows that $mx^{m-1}t$ is not divisible by $p$. This implies that $y^m\nequiv x^m\pmod{p^{s+1}}$. On the other hand, the congruence $x\equiv y\pmod{p^s}$ clearly implies that $x^m\equiv y^m\pmod{p^s}$, so we conclude that
  \begin{equation*}
    \nu_p(x^m-y^m)=s.
  \end{equation*}
  The inductive step is a direct consequence of \cref{lemma:lte-step}.
\end{proof}

\begin{definition}
  A strong divisibility sequence is said to satisfy the \emph{law of repetition} if for any prime $p$, we have
  \begin{equation*}
    \nu_p(a_{n\rho_a(p)})=\nu_p(a_{\rho_a(p)})+\nu_p(n).
  \end{equation*}
\end{definition}

\begin{theorem}
  Consider a strong divisibility sequence $(a_n)$, such that for each prime $p$ we have
  \begin{equation*}
    \rho_a(p^k)=\rho_a(p)p^{k-1}.
  \end{equation*}
  Then $(a_n)$ satisfies the law of repetition.
\end{theorem}

\begin{example}
  For any two relatively prime integers $a$ and $b$, the sequence
  \begin{equation*}
    (a^n-b^n)
  \end{equation*}
  satisfies the law of repetition.
\end{example}

\begin{example}
  Any Lucas sequence satisfies the law of repetition.
\end{example}

\exercises

\begin{routine}
    \exercise Show that
  \begin{equation*}
    \gcd(a^m-b^m,a^n-b^n)=a^{\gcd(m,n)}-b^{\gcd(m,n)}.
  \end{equation*}
  \exercise Consider a positive integer $a$, and write $S_m:=\sum_{k=0}^m a^k$. Prove that
  \begin{equation*}
    \gcd(S_m,S_n)=S_{\gcd(m+1,n+1)-1}.
  \end{equation*}
\end{routine}

\begin{challenge}
    \exercise
  \begin{subexenum}
  \item Show that the function $(x,y)\mapsto (y,x+y)$ is a bijection on the set of pairs of elements $x,y\in \mathbb{Z}/n\mathbb{Z}$.
  \item Show that for any $m$, the Fibonacci sequence is periodic modulo $m$. Its period is called the \emph{Pisano period}.
  \item Show that for any natural number $m$, there is a least positive number $n$ such that
    \begin{equation*}
      m\mid F_n.
    \end{equation*}
    We will call this number $G_m$. By virtue of its duality to the Fibonacci sequence, we will call the sequence $G$ the \emph{cofibonacci sequence}.
  \item Show that $G(m)$ satisfies the following logical equivalence
    \begin{equation*}
      G_m\mid n \Leftrightarrow m \mid F_n.
    \end{equation*}
    Conclude that the Fibonacci sequence is a strong divisibility sequence.
  \end{subexenum}
  \exercise Use methods similar to the previous exercise to show that the sequence of Pell numbers $P_n$, which are defined by
  \begin{equation*}
    P_0:=0,\qquad P_1:=1,\qquad\text{and}\qquad P_{n+2}=2P_{n+1}+P_{n},
  \end{equation*}
  forms a strong divisibility sequence.
\end{challenge}

%%% Local Variables:
%%% mode: latex
%%% TeX-master: "elementary-number-theory"
%%% TeX-engine: luatex
%%% End:
